\section{Conclusion and Future Work}\label{sec:conclusion}

We have shown that clairvoyance and demand semantics are two sides of the same
coin using logic programming as an intermediary. This connection is both
theoretically interesting and practically useful: it demonstrates that
functional logic programming concepts have value in seemingly unrelated
areas. We also obtained an executable clairvoyance semantics and a simpler demand
semantics built on functional logic programming semantics.

As discussed in \cref{sec:case-study}, our technique is slow with even
modestly-sized inputs. Although unfortunate, this may be acceptable for most applications. A major
motivation behind this line of research was the desire to find a lightweight
representation of demand functions that could be used to test their asymptotic
bounds (constants included) \emph{before} trying to formally verify them. For
this purpose, it usually suffices to test quite small inputs, where our
technique performs quite well in absolute (wall-clock) terms. However, some
applications seem to perform much worse than others for reasons that we do not
yet understand. Further investigation is left as future work.
